% 地支三会图
% 用途：展示四个三会局，连续三个地支成一方之气
\documentclass[tikz,border=10pt]{standalone}
\usepackage{ctex}
\usepackage{xcolor}
\pagecolor{white}
\usetikzlibrary{arrows.meta, positioning, calc, decorations.pathmorphing, backgrounds}

\begin{document}
\begin{tikzpicture}[scale=1.3]
  \coordinate (zi)  at ( 90:3.8);
  \coordinate (chou) at ( 60:3.8);
  \coordinate (yin)  at ( 30:3.8);
  \coordinate (mao)  at (  0:3.8);
  \coordinate (chen) at (330:3.8);
  \coordinate (si)   at (300:3.8);
  \coordinate (wu)   at (270:3.8);
  \coordinate (wei)  at (240:3.8);
  \coordinate (shen) at (210:3.8);
  \coordinate (you)  at (180:3.8);
  \coordinate (xu)   at (150:3.8);
  \coordinate (hai)  at (120:3.8);

  % 三会扇形背景
  \begin{scope}[on background layer]
    \fill[green!15]  (30:4.1) arc(30:150:4.1) -- (0,0) -- cycle;
    \fill[red!15]    (330:4.1) arc(330:30:4.1) -- (0,0) -- cycle;
    \fill[orange!15] (210:4.1) arc(210:330:4.1) -- (0,0) -- cycle;
    \fill[blue!15]   (150:4.1) arc(150:210:4.1) -- (0,0) -- cycle;
  \end{scope}

  \draw[thick, gray!60] (0,0) circle (3.8);

  \foreach \name/\angle/\text in {
    zi/90/子, chou/60/丑, yin/30/寅, mao/0/卯,
    chen/330/辰, si/300/巳, wu/270/午, wei/240/未,
    shen/210/申, you/180/酉, xu/150/戌, hai/120/亥
  } {
    \fill[black] (\angle:3.8) circle (2pt);
    \node[font=\Large\bfseries] at (\angle:4.6) {\text};
  }

  % 三会局标签，分别放在四个方向区域
  % 北方水：上方
  \node[font=\normalsize\bfseries, blue!60!black] at (90:2.5) {北方水};
  \node[font=\normalsize\bfseries, blue!60!black] at (90:1.8) {亥子丑};
  % 东方木：右侧
  \node[font=\normalsize\bfseries, green!60!black] at (0:2.5) {东方木};
  \node[font=\normalsize\bfseries, green!60!black] at (0:1.8) {寅卯辰};
  % 南方火：下方
  \node[font=\normalsize\bfseries, red!60!black] at (270:2.5) {南方火};
  \node[font=\normalsize\bfseries, red!60!black] at (270:1.8) {巳午未};
  % 西方金：左侧
  \node[font=\normalsize\bfseries, orange!60!black] at (180:2.5) {西方金};
  \node[font=\normalsize\bfseries, orange!60!black] at (180:1.8) {申酉戌};

  % 图例
  \node[font=\small, anchor=west] at (-5.8,-5.2) {三会：三个连续地支连成一方之气};

  % 白色背景，防止 GitHub dark mode 下透明 SVG 看不清
  \begin{scope}[on background layer]
    \fill[white] (current bounding box.south west) rectangle (current bounding box.north east);
  \end{scope}
\end{tikzpicture}
\end{document}
